\section{RQ3: Quantum-Vulnerable Cryptography in Critical Open Source}
\label{sec:rq3}

\subsection{Method}

\paragraph{Projects.} The OpenSSF Criticality Score ranks open-source
projects by signals of how much depends on them~\cite{criticality}. We used its
latest published snapshot, \data{rqthree-snapshot}, pinned by commit and hash,
and walked down the top \data{rqthree-candidates} GitHub repositories. Many
top entries are copies of the same code base, mostly Linux kernel trees, so a
repository whose root commit matches a higher-ranked one was skipped as a
duplicate: \data{rqthree-duplicates} were, \data{rqthree-duplicates-kernel} of
them kernel trees. The first \data{rqthree-projects} distinct code bases, down
to rank \data{rqthree-last-rank}, are the sample. Repositories above
\data{rqthree-cap-gb}\,GB were counted but not cloned:
\data{rqthree-skipped-list}. One clone, \data{rqthree-failed-list}, failed
twice. That leaves \data{rqthree-scanned} scanned projects.

\paragraph{Scan.} Each project was cloned shallow at its default branch's head
and scanned with \tool{}'s source, configuration, dependency and binary layers
and reconciliation, with a \data{rqthree-timeout-min}-minute limit that no scan
reached (median \data{rqthree-scan-median}). Nothing was built or run. A
finding counts when its asset is quantum-vulnerable and it was found by a
layer rather than derived; derived findings are reported separately. Each
finding's exposure class comes from the risk engine: \emph{test} for test and
example paths; otherwise \emph{network-facing} for configuration and
key-material findings and for calls into TLS or SSH APIs; and
\emph{internal} for the rest.

\subsection{Results}

\begin{table}[t]
\centering\footnotesize\setlength{\tabcolsep}{4pt}
% Generated by paper/gen.py from paper/data. Do not edit.
\begin{tabular}{lrrrr}
\toprule
Projects with a finding & $n$ & Share & 95\% CI & Findings \\
\midrule
Any layer & 50 & 51.5\% & 42--61 & 7,596 \\
Used (source, config, binary) & 48 & 49.5\% & 40--59 & 6,323 \\
\emph{By layer} &  &  &  &  \\
\quad source & 32 & 33.0\% & 24--43 & 671 \\
\quad config and keys & 41 & 42.3\% & 33--52 & 5,637 \\
\quad dependencies & 26 & 26.8\% & 19--36 & 1,273 \\
\quad binaries & 2 & 2.1\% & 1--7 & 15 \\
\emph{By exposure class} &  &  &  &  \\
\quad network-facing & 22 & 22.7\% & 15--32 & 1,596 \\
\quad internal & 34 & 35.0\% & 26--45 & 1,366 \\
\quad test and example & 41 & 42.3\% & 33--52 & 4,634 \\
\emph{Derived findings} &  &  &  &  \\
\quad binary without source & 2 & 2.1\% & 1--7 & 9 \\
\quad stale dependency & 7 & 7.2\% & 4--14 & 20 \\
\bottomrule
\end{tabular}

\caption{Scanned projects with at least one quantum-vulnerable finding
($n$ of \data{rqthree-scanned}), with Wilson 95\% intervals, and the number
of findings. ``Used'' excludes dependency findings, which show a capability,
not a use.}
\label{tab:rq3-summary}
\end{table}

\paragraph{Prevalence.} Table~\ref{tab:rq3-summary} summarizes the scan.
\data{rqthree-any-n} of \data{rqthree-scanned} projects
(\data{rqthree-any-share}, CI \data{rqthree-any-ci}) contain
quantum-vulnerable cryptography in some layer, and
\data{rqthree-used-n} (\data{rqthree-used-share}) in their own source,
configuration or binaries, not only in dependencies. Configuration and key
material is the most common layer (\data{rqthree-config-share}), ahead of
source code (\data{rqthree-source-share}) and dependencies
(\data{rqthree-deps-share}); vulnerable code in committed binaries is rare
(\data{rqthree-binary-share}).

\paragraph{Exposure.} \data{rqthree-network-n} projects
(\data{rqthree-network-share}, CI \data{rqthree-network-ci}) have
quantum-vulnerable cryptography in network-facing code or configuration.
Test and example code is more common: \data{rqthree-test-share} of projects
have quantum-vulnerable findings there.

\paragraph{Algorithms.} RSA is the most widespread algorithm, in
\data{rqthree-alg-rsa-share} of scanned projects, followed by ECDSA
(\data{rqthree-alg-ecdsa-share}) and EdDSA (\data{rqthree-alg-eddsa-share});
Table~\ref{tab:rq3-algorithms} in the appendix lists all.

\paragraph{Reconciliation.} \data{rqthree-derived-stale-n} projects pin a
library version too old to offer post-quantum algorithms, and
\data{rqthree-derived-bws-n} ship binaries with cryptography their source
scan does not show. The other two derived types did not occur.
Declared-PQC-deployed-classical needs a probed endpoint, and RQ3 probes none.

\paragraph{Counts are concentrated.} Findings are not a measure of
prevalence. Five projects (\data{rqthree-top5-list}) account for
\data{rqthree-top5-share} of all quantum-vulnerable findings,
\data{rqthree-top5-config-share} of them certificates, keys and
configuration. The median
project with any finding has \data{rqthree-median-findings}. We report
shares of projects for this reason.

\subsection{Ranking Findings for Migration}

A project with hundreds of findings needs an order to migrate them in, which
is what \tool{}'s risk score gives (\S\ref{sec:design}). Its four factors
(exposure, secrecy lifetime, quantum status and confidence) carry weights we
chose by judgment, so we ablate them on the RQ3 findings. Each factor in turn
is set to one for every finding, the findings are re-scored, and each
project's ranking is compared with the full score's. That covers the
\data{riskabl-findings} findings of the \data{riskabl-projects} projects with
at least two, and the \data{riskabl-pairs} pairs of findings within a
project that the full score orders.

\begin{table}[t]
\centering\footnotesize\setlength{\tabcolsep}{4pt}
% Generated by paper/gen.py from paper/data. Do not edit.
\begin{tabular}{lrrrr}
\toprule
Factor removed & Reversed & Tied & Top changed & Median $\tau_b$ \\
\midrule
exposure & 6.8\% & 9.3\% & 45 (62\%) & 0.786 \\
secrecy lifetime & 12.1\% & 1.4\% & 28 (39\%) & 0.897 \\
quantum status & 16.4\% & 4.4\% & 34 (47\%) & 0.593 \\
confidence & 0.8\% & 1.2\% & 19 (26\%) & 0.992 \\
\bottomrule
\end{tabular}

\caption{Risk-score ablation over the \data{riskabl-projects} RQ3 projects
with at least two findings. Reversed and tied: share of the pairs the full
score orders that the ablated score orders the other way or ties. Top set
changed: projects whose highest-risk findings differ. $\tau_b$: Kendall's
rank correlation with the full score, median over projects.}
\label{tab:risk-ablation}
\end{table}

Table~\ref{tab:risk-ablation} gives the results. Every factor changes the
ranking, in different ways. Quantum status shapes the overall order most:
without it \data{riskabl-quantum-status-reversed} of ordered pairs reverse
and the median $\tau_b$ falls to \data{riskabl-quantum-status-tau}, because
hashes and symmetric ciphers then rank with RSA and elliptic curves.
Exposure decides what comes first: without it the top-risk findings change in
\data{riskabl-exposure-top} of \data{riskabl-projects} projects
(\data{riskabl-exposure-top-share}), and \data{riskabl-exposure-tied} of
pairs become ties. Secrecy lifetime reverses
\data{riskabl-secrecy-lifetime-reversed} of pairs. Confidence matters least: \data{riskabl-confidence-reversed}
of pairs reverse and the median $\tau_b$ stays at
\data{riskabl-confidence-tau}. It still changes the top set in \data{riskabl-confidence-top} projects. The
pair counts are pooled, so projects with many findings weigh most; the
$\tau_b$ median and the top-set count weigh each project once. The ablation
shows that each factor matters to the order, not that its weights are right:
validating them would need migration priorities set by the projects
themselves.

\paragraph{Limits.} \tool{}'s rules match cryptographic APIs, libraries and
configuration. A project that implements an algorithm itself shows few
findings: the Linux kernel tree in the sample, \data{rqthree-kernel-repo},
has \data{rqthree-kernel-findings}. Source rules cover Go, Python, Java, C/C++
and JavaScript/TypeScript; \data{rqthree-uncovered} scanned projects have a
primary language outside them (\data{rqthree-uncovered-top} and
\data{rqthree-uncovered-others} others) and are scanned in
the other layers only. Among the \data{rqthree-covered} in covered languages,
\data{rqthree-covered-source-share} have vulnerable source findings. The
network-facing class is a path and API heuristic, not reachability: it misses
cryptography reached from a network handler through other calls and includes
TLS configuration that no deployment uses. Finally, one shallow commit per
project is scanned on a case-insensitive file system, which drops one path of
each pair differing only in case.
